    
    
    

    
    \hypertarget{newton-raphson-method-for-finding-root}{%
\section{Newton-Raphson method for finding
root}\label{newton-raphson-method-for-finding-root}}

    Newton-Raphson method is a root-finding algorithm to find the root
(i.e., the zeroes) of a real-valued function using numerical analysis.
The Newton-Raphson method produces succesive better approximations of
the root with each iteration of any real function.

For finding the root of a single-variable real-valued function \(f(x)\)
the algorithm requires the function itself, the derivative of the
function \(f'(x)\), and an initial guess \(x_0\) for the root of \(f\).

\begin{figure}
\centering
\includegraphics{resources/newton-raphson-1.png}
\caption{resources/newton-raphson-1.png}
\end{figure}

Using \(f'(x)\) we first evaluate the tangent to the curve \(f(x)\) at
\(x = x_0\). Then, our next guess \(x_1\), becomes the intersection of
this tangent at the \(x-axis\). From the diagram above, the derivetive
of \(f(x)\) at \(x = x_0\) will be given by,
\[f'(x_0) = \frac{f(x_0) - 0}{x_0 - x_1}\]

Thus, the initial guess \(x_1\) would be given by,
\[x_1 = x_0 - \frac{f(x_0)}{f'(x_0)}\]

The next guess \(x_2\) would be, \[x_2 = x_1 - \frac{f(x_1)}{f'(x_1)}\]

This process is to be repeated again and again until a guess close
enough to the actual root is attained.

Therefore, the generic form of the \(nth\) guess for the root would be,
\[x_n = x_{n-1} - \frac{f(x_{n-1})}{f'(x_{n-1})}\]

    \hypertarget{code}{%
\subsection{Code}\label{code}}

    \begin{tcolorbox}[breakable, size=fbox, boxrule=1pt, pad at break*=1mm,colback=cellbackground, colframe=cellborder]
\prompt{In}{incolor}{1}{\boxspacing}
\begin{Verbatim}[commandchars=\\\{\}]
\PY{k+kn}{from} \PY{n+nn}{numpy} \PY{k+kn}{import} \PY{n}{cos}\PY{p}{,} \PY{n}{e}\PY{p}{,} \PY{n}{linspace}\PY{p}{,} \PY{n}{pi}\PY{p}{,} \PY{n}{sin}
\end{Verbatim}
\end{tcolorbox}

    \begin{tcolorbox}[breakable, size=fbox, boxrule=1pt, pad at break*=1mm,colback=cellbackground, colframe=cellborder]
\prompt{In}{incolor}{2}{\boxspacing}
\begin{Verbatim}[commandchars=\\\{\}]
\PY{c+c1}{\PYZsh{} Code for finding root using Newton\PYZhy{}Raphson method}
\PY{k}{def} \PY{n+nf}{find\PYZus{}root\PYZus{}newton}\PY{p}{(}\PY{n}{f}\PY{p}{,} \PY{n}{f\PYZus{}prime}\PY{p}{,} \PY{n}{x\PYZus{}0}\PY{p}{,} \PY{n}{tolerance}\PY{o}{=}\PY{l+m+mf}{0.001}\PY{p}{)}\PY{p}{:}

    \PY{c+c1}{\PYZsh{} Declaring a variable to store the `f(x) / f\PYZus{}prime(x)`}
    \PY{n}{h} \PY{o}{=} \PY{p}{(}\PY{n}{f}\PY{p}{(}\PY{n}{x\PYZus{}0}\PY{p}{)} \PY{o}{/} \PY{n}{f\PYZus{}prime}\PY{p}{(}\PY{n}{x\PYZus{}0}\PY{p}{)}\PY{p}{)}

    \PY{c+c1}{\PYZsh{} Evaluating the root}
    \PY{k}{while} \PY{n+nb}{abs}\PY{p}{(}\PY{n}{h}\PY{p}{)} \PY{o}{\PYZgt{}} \PY{n}{tolerance}\PY{p}{:}

        \PY{c+c1}{\PYZsh{} Evaluating `f(x) / f\PYZus{}prime(x)`}
        \PY{n}{h} \PY{o}{=} \PY{p}{(}\PY{n}{f}\PY{p}{(}\PY{n}{x\PYZus{}0}\PY{p}{)} \PY{o}{/} \PY{n}{f\PYZus{}prime}\PY{p}{(}\PY{n}{x\PYZus{}0}\PY{p}{)}\PY{p}{)}

        \PY{c+c1}{\PYZsh{} Subtracting `f(x) / f\PYZus{}prime(x)` from the previous guess}
        \PY{n}{x\PYZus{}0} \PY{o}{\PYZhy{}}\PY{o}{=} \PY{n}{h}

    \PY{k}{return} \PY{n}{x\PYZus{}0}
\end{Verbatim}
\end{tcolorbox}

    \hypertarget{examples}{%
\subsection{Examples}\label{examples}}

    \hypertarget{x2-2}{%
\subsubsection{\texorpdfstring{\[x^2 = 2\]}{x\^{}2 = 2}}\label{x2-2}}

\[f(x) = x^2 - 2\] \[f'(x) = 2x\]

    \begin{tcolorbox}[breakable, size=fbox, boxrule=1pt, pad at break*=1mm,colback=cellbackground, colframe=cellborder]
\prompt{In}{incolor}{3}{\boxspacing}
\begin{Verbatim}[commandchars=\\\{\}]
\PY{k}{def} \PY{n+nf}{f}\PY{p}{(}\PY{n}{x}\PY{p}{)}\PY{p}{:}
    \PY{k}{return} \PY{n}{x}\PY{o}{*}\PY{o}{*}\PY{l+m+mi}{2} \PY{o}{\PYZhy{}} \PY{l+m+mi}{2}


\PY{k}{def} \PY{n+nf}{f\PYZus{}prime}\PY{p}{(}\PY{n}{x}\PY{p}{)}\PY{p}{:}
    \PY{k}{return} \PY{l+m+mi}{2} \PY{o}{*} \PY{n}{x}


\PY{n}{x\PYZus{}0} \PY{o}{=} \PY{n}{pi}

\PY{n}{find\PYZus{}root\PYZus{}newton}\PY{p}{(}\PY{n}{f}\PY{p}{,} \PY{n}{f\PYZus{}prime}\PY{p}{,} \PY{n}{x\PYZus{}0}\PY{p}{)}
\end{Verbatim}
\end{tcolorbox}

            \begin{tcolorbox}[breakable, size=fbox, boxrule=.5pt, pad at break*=1mm, opacityfill=0]
\prompt{Out}{outcolor}{3}{\boxspacing}
\begin{Verbatim}[commandchars=\\\{\}]
1.414213562373189
\end{Verbatim}
\end{tcolorbox}
        
    \hypertarget{x3---x2-x-e}{%
\subsubsection{\texorpdfstring{\[x^3 - x^2 + x = e\]}{x\^{}3 - x\^{}2 + x = e}}\label{x3---x2-x-e}}

\[f(x) = x^3 - x^2 + x - e\] \[f'(x) = 3x^2 - 2x + 1\]

    \begin{tcolorbox}[breakable, size=fbox, boxrule=1pt, pad at break*=1mm,colback=cellbackground, colframe=cellborder]
\prompt{In}{incolor}{4}{\boxspacing}
\begin{Verbatim}[commandchars=\\\{\}]
\PY{k}{def} \PY{n+nf}{f}\PY{p}{(}\PY{n}{x}\PY{p}{)}\PY{p}{:}
    \PY{k}{return} \PY{n}{x}\PY{o}{*}\PY{o}{*}\PY{l+m+mi}{3} \PY{o}{\PYZhy{}} \PY{n}{x}\PY{o}{*}\PY{o}{*}\PY{l+m+mi}{2} \PY{o}{+} \PY{n}{x} \PY{o}{\PYZhy{}} \PY{n}{e}


\PY{k}{def} \PY{n+nf}{f\PYZus{}prime}\PY{p}{(}\PY{n}{x}\PY{p}{)}\PY{p}{:}
    \PY{k}{return} \PY{l+m+mi}{3} \PY{o}{*} \PY{n}{x}\PY{o}{*}\PY{o}{*}\PY{l+m+mi}{2} \PY{o}{\PYZhy{}} \PY{l+m+mi}{2} \PY{o}{*} \PY{n}{x} \PY{o}{+} \PY{l+m+mi}{1}


\PY{n}{x\PYZus{}0} \PY{o}{=} \PY{l+m+mi}{0}

\PY{n}{find\PYZus{}root\PYZus{}newton}\PY{p}{(}\PY{n}{f}\PY{p}{,} \PY{n}{f\PYZus{}prime}\PY{p}{,} \PY{n}{x\PYZus{}0}\PY{p}{)}
\end{Verbatim}
\end{tcolorbox}

            \begin{tcolorbox}[breakable, size=fbox, boxrule=.5pt, pad at break*=1mm, opacityfill=0]
\prompt{Out}{outcolor}{4}{\boxspacing}
\begin{Verbatim}[commandchars=\\\{\}]
1.5193605629027016
\end{Verbatim}
\end{tcolorbox}
        

    % Add a bibliography block to the postdoc
    
    
    
